\section{联合数据--模型 Scaling Laws 与 Chinchilla}

\subsection{联合 Scaling 的动机}

前面我们讨论的 Scaling Laws 都是\textbf{单变量}的：固定其他变量，只看一个变量（数据、参数或计算）对损失的影响。但在实际训练中，\textbf{计算预算是固定的}，你需要决定如何在模型大小和数据量之间分配。

\begin{figure}[H]
\centering
\includegraphics[width=0.95\textwidth]{slides-images/slide-032.jpg}
\caption{核心问题：给定固定计算预算，模型应该多大？数据应该多少？\protect\footnotemark}
\end{figure}
\footnotetext{Slides 第33页。}

两个极端都不好：
\begin{itemize}
    \item 极小模型 + 海量数据 $\rightarrow$ 模型容量不足，数据被浪费
    \item 极大模型 + 极少数据 $\rightarrow$ 严重过拟合，计算被浪费
\end{itemize}

\subsection{联合 Scaling Law 的函数形式}

\begin{figure}[H]
\centering
\includegraphics[width=0.95\textwidth]{slides-images/slide-033.jpg}
\caption{Rosenfeld 和 Kaplan 的联合 Scaling Law 函数形式\protect\footnotemark}
\end{figure}
\footnotetext{Slides 第34页。}

Rosenfeld et al. 提出的函数形式：
$$L(N, D) = \left(\frac{N_c}{N}\right)^{\alpha_N} + \left(\frac{D_c}{D}\right)^{\alpha_D} + \epsilon_\infty$$

其中 $N$ 是参数量，$D$ 是数据量（token 数），$\epsilon_\infty$ 是不可约误差。每项分别描述了模型容量不足（$N$ 项）和数据不足（$D$ 项）带来的误差。

Kaplan et al. 使用了类似但稍有不同的形式（聚焦于可约误差而非不可约误差）。

\begin{figure}[H]
\centering
\includegraphics[width=0.95\textwidth]{slides-images/slide-034.jpg}
\caption{Rosenfeld 的 3D Scaling Law 拟合：数据量 $\times$ 模型大小 $\times$ 损失\protect\footnotemark}
\end{figure}
\footnotetext{Slides 第35页。}

\begin{importantbox}{令人惊讶的拟合精度}
尽管联合 Scaling Law 的函数形式是``凭经验拍脑袋''（ad hoc）的，但它对实际训练结果的拟合\textbf{几乎完美}。Rosenfeld 展示了只用小模型+少数据拟合的 Scaling Law，可以非常准确地外推到大模型+多数据的区域。
\end{importantbox}

\begin{figure}[H]
\centering
\includegraphics[width=0.95\textwidth]{slides-images/slide-035.jpg}
\caption{Rosenfeld：用小规模训练结果准确外推大规模性能\protect\footnotemark}
\end{figure}
\footnotetext{Slides 第36页。}

\subsection{Kaplan 的固定计算预算分析}

\begin{figure}[H]
\centering
\includegraphics[width=0.95\textwidth]{slides-images/slide-036.jpg}
\caption{Kaplan：固定计算量下，参数量与损失的关系（不同颜色代表不同计算量）\protect\footnotemark}
\end{figure}
\footnotetext{Slides 第37页。}

在 Kaplan 的论文中，固定计算量（颜色编码），变化参数量（x轴），可以看到每条等计算曲线上都有一个最优点——对应着在该计算预算下的最优模型大小。

然而，Kaplan 的最优估计后来被发现有较大偏差。原因之一是学习率 schedule 的影响（cosine schedule 不能被提前截断），导致在不同计算预算下比较模型时引入了系统性偏差。

\subsection{Chinchilla：三种方法}

Chinchilla 论文（Hoffmann et al., Google DeepMind）是确定最优数据/模型比例的权威参考。

\begin{figure}[H]
\centering
\includegraphics[width=0.95\textwidth]{slides-images/slide-038.jpg}
\caption{Chinchilla 三种方法的估计以及与 Kaplan 的对比\protect\footnotemark}
\end{figure}
\footnotetext{Slides 第39页。}

\begin{table}[H]
\centering
\begin{tabular}{lcc}
\toprule
\textbf{方法} & \textbf{参数 Scaling 指数 $a$} & \textbf{数据 Scaling 指数 $b$} \\
\midrule
Chinchilla 方法 1 & 0.50 & 0.50 \\
Chinchilla 方法 2 & 0.50 & 0.50 \\
Chinchilla 方法 3 & 0.46 & 0.54 \\
Kaplan et al. & 0.73 & 0.27 \\
\bottomrule
\end{tabular}
\caption{不同方法估计的 Scaling 指数：$N_{\text{opt}} \propto C^a$, $D_{\text{opt}} \propto C^b$}
\end{table}

\textbf{方法 1：最小包络法}

\begin{figure}[H]
\centering
\includegraphics[width=0.95\textwidth]{slides-images/slide-039.jpg}
\caption{Chinchilla 方法 1：取训练曲线的下包络，找出每个计算预算下的最优模型\protect\footnotemark}
\end{figure}
\footnotetext{Slides 第40页。}

在所有不同大小的模型训练曲线上，取\textbf{下包络线}（lower envelope）——即对于每个 FLOP 值，选择达到最低损失的那个模型。然后分析这些最优点的参数量和 token 数如何随 FLOP 缩放。

\textbf{方法 2：IsoFLOP 分析}

\begin{figure}[H]
\centering
\includegraphics[width=0.95\textwidth]{slides-images/slide-041.jpg}
\caption{Chinchilla 方法 2（IsoFLOP）：固定计算量，扫描模型大小，找最优点\protect\footnotemark}
\end{figure}
\footnotetext{Slides 第42页。}

\begin{importantbox}{IsoFLOP 分析——最规范的方法}
IsoFLOP 分析可能是最直观、最规范的 Chinchilla 方法：
\begin{enumerate}
    \item 选择多个计算尺度（每种颜色代表一个固定的 FLOP 预算）
    \item 在每个尺度上，训练不同大小的模型（小模型+多数据 到 大模型+少数据）
    \item 找到每条等 FLOP 曲线的最小值点
    \item 分析最小值点如何随 FLOP 缩放 $\rightarrow$ 得到最优参数量和最优 token 数的缩放关系
\end{enumerate}
\end{importantbox}

\begin{figure}[H]
\centering
\includegraphics[width=0.95\textwidth]{slides-images/slide-042.jpg}
\caption{IsoFLOP 分析的结果：最优参数量和最优 token 数与 FLOP 的关系\protect\footnotemark}
\end{figure}
\footnotetext{Slides 第43页。}

方法 1 和方法 2 给出了高度一致的结果：对于 Chinchilla 的训练预算，最优模型约 63--67B 参数。

\textbf{方法 3：参数化曲面拟合}

\begin{figure}[H]
\centering
\includegraphics[width=0.95\textwidth]{slides-images/slide-043.jpg}
\caption{Chinchilla 方法 3：拟合参数化 Scaling Law 曲面\protect\footnotemark}
\end{figure}
\footnotetext{Slides 第44页。}

方法 3 是直接拟合 Rosenfeld 式的联合 Scaling Law 函数。训练大量模型（不同参数量 $\times$ 不同数据量），拟合 3D 曲面，然后求最优。但这个方法的拟合质量明显不如前两种。

\begin{warningbox}{Chinchilla 方法 3 的曲线拟合 bug}
Epoch AI 的研究者在 2024 年试图复现 Chinchilla 方法 3 的结果。由于无法获取原始数据，他们甚至用``图表取值器''从论文图表中提取数据点。结果发现：

\begin{itemize}
    \item 原始论文的残差\textbf{不是零均值}的——这是曲线拟合有问题的标志
    \item 修正拟合后，方法 3 的估计与方法 1、2 几乎完全一致
    \item 即原始结果是正确的，只是拟合过程中有小错误
\end{itemize}

这是一个``复现反而证实了原始结论''的罕见案例。
\end{warningbox}

\subsection{Chinchilla 比例与 Cosine Schedule 的陷阱}

\begin{figure}[H]
\centering
\includegraphics[width=0.95\textwidth]{slides-images/slide-044.jpg}
\caption{Cosine 学习率 Schedule 的陷阱\protect\footnotemark}
\end{figure}
\footnotetext{Slides 第45页。}

Kaplan 估计偏差较大的一个重要原因是 \textbf{cosine learning rate schedule 不能提前截断}：

\begin{warningbox}{Cosine Schedule 的不可截断性}
Cosine 学习率先升后降，必须完整运行到冷却（cool-down）阶段才能得到有效的模型。如果在中途截断，模型仍处于高学习率状态，性能远差于完整训练的模型。因此，\textbf{在做 Scaling Law 实验时，每个模型必须用完整的 schedule 训练到底}，不能简单地取训练中间的 checkpoint 来做对比。
\end{warningbox}

著名的 Chinchilla 比例约为 \textbf{20 tokens/parameter}。这意味着一个 70B 参数的模型最优应训练 1.4T tokens。

\subsection{本章小结}

联合数据--模型 Scaling Laws 解答了大语言模型训练中最核心的资源分配问题。Chinchilla 的三种分析方法从不同角度收敛到了相似的结论（$a \approx b \approx 0.5$），确立了约 20 tokens/parameter 的黄金比例。

%% ========== Part 8: 训练最优 vs. 推理最优 ==========

